\thispagestyle{empty}
{\sffamily\small PROVEIT\quad SOURCE REVIEW\par}
\vspace{10pt}
{\fontsize{25}{29}\selectfont\bfseries\raggedright Simplifying Lean proofs\\in DiophantineConsistency\par}
\vspace{8pt}
{\large Source specific refactoring proposals\par}
\vspace{7pt}
26 October 2026\qquad Pinned static review
\vspace{15pt}

\textbf{Start with reuse, not stronger automation.} The clearest first edit
is \code{half_base_bound_31} in \code{SelectiveCoefficientBounds}: it repeats the
generic theorem immediately above it. Specializing that existing result at
31 produces a two-line candidate for the current 24-line proof body, with
the same strict conclusion and assumptions. Next remove exact duplicated
coordinate-packaging arguments and give long conjunction certificates named
accessors. These changes make the intended mathematics more visible.

For the large generated proofs, the necessary DAG soundness, power, and
degree APIs already exist. The useful next step is to route generated
instances through them, retaining the literal circuit and public theorem
contracts. The supplied conservative experiment replaces 193 individual
power \code{ring} proofs with shared finite-label certificates while keeping
all 1,192 public gate/power theorem header strings unchanged.

\textbf{The measured circuit tradeoff is honest rather than spectacular:}
660 fewer physical lines, but 2,360 more UTF-8 bytes in the final candidates.
This is less repeated proof scripting, not byte compression, a smaller
circuit, or a demonstrated speedup. More compact indexed facades are
included as experiments, not advertised as drop-in replacements.

\textbf{Status and trust.} No Lean or Lake executable was available.
\textbf{Every proposed Lean candidate is uncompiled.} This report supplies
source-specific designs, exact pinned references, practical before/after
edits, complete candidate artifacts, static evidence, and a replay roadmap.
It does not claim an accepted patch, kernel verification, faster builds,
or smaller elaborated proofs. The original source snapshot is unchanged.

\textbf{Mathematical scope.} Preserve the project's standard-axiom policy
and every theorem's witness domain and hypotheses. This is not a proof of
\(\operatorname{Con}(\mathrm{ZFC})\), a consistency decision, or completion
of the open polynomial/native-machine or internal-ZFC links identified in
\docsrc{README.md}{14}{31}.

\vspace{8pt}
\textbf{Recommended first patch.} Try the generic-bound specialization
(Section~\ref{sec:halfbase}), the finite-set support identities, and the
sum-of-squares proof (Section~\ref{sec:squares}). Replay them against the
pinned baseline. Then compile all four compatibility tests for the shared
nine-coordinate helper (Section~\ref{sec:nine}). Leave representation
changes and global automation until these small pilots establish the
validation workflow.

\newpage
\section*{Recommendation map}

Priority 1 means a small interface-preserving pilot. Priority 2 is the next
maintenance/refactoring batch. Priority 3 needs a measured consumer or
abstraction experiment. Priority 4 is a later design option, not an immediate
rewrite recommendation.

\small
\begin{tabularx}{\linewidth}{@{}L{0.42in}YL{1.35in}L{0.56in}@{}}
\toprule
Rank & Change & Main caution & Section\\
\midrule
1 & Specialize the existing 31-variable coefficient bound & Compile the exact statement & \ref{sec:halfbase}\\[4pt]
1 & Use sum-of-squares and finite-set identities & Conditional simp needs replay & \ref{sec:local}\\[4pt]
1 & Share nine-coordinate packaging and add named projections & Preserve slot 7 and positivity & \ref{sec:nine}, \ref{sec:accessors}\\[4pt]
2 & Reuse \code{run_power}, keep public wrappers & Certificate cost is unmeasured & \ref{sec:powercert}\\[4pt]
2 & Share source/window and modular-residue interfaces & Same source and parameters & \ref{sec:window}, \ref{sec:residue}\\[4pt]
2 & Generate pure audit wrappers and enforce axiom policy & Keep behavioral checks & \ref{sec:axiomgate}\\[4pt]
3 & Pilot indexed gates, shared degree profiles, and degree prefixes & Keep semantic cut points & \ref{sec:indexed}, \ref{sec:profiles}, \ref{sec:shortdegree}\\[4pt]
4 & Unify coefficient syntax or add typed circuit views & Representation bridges required & \ref{sec:coefficients}, \ref{sec:circuits}\\
\bottomrule
\end{tabularx}
\normalsize

The archive's before/after comparisons are measured text spans, not additive
project-wide forecasts. In particular, the predecessor-exponent cleanup
overlaps the 31-variable specialization, and the concise/explicit
sum-of-squares proofs are alternatives. The 60-line tuple duplicate is an
opportunity for shared maintenance, not a guaranteed net deletion after
adding a helper and compatibility wrappers.

\vspace{5pt}
{\small\setlength{\parskip}{1pt}\tableofcontents}
